% Options for packages loaded elsewhere
\PassOptionsToPackage{unicode}{hyperref}
\PassOptionsToPackage{hyphens}{url}
%
\documentclass[
  ignorenonframetext,
  aspectratio=169,
]{beamer}
\usepackage{pgfpages}
\setbeamertemplate{caption}[numbered]
\setbeamertemplate{caption label separator}{: }
\setbeamercolor{caption name}{fg=normal text.fg}
\beamertemplatenavigationsymbolshorizontal
% Prevent slide breaks in the middle of a paragraph
\widowpenalties 1 10000
\raggedbottom
\setbeamertemplate{part page}{
  \centering
  \begin{beamercolorbox}[sep=16pt,center]{part title}
    \usebeamerfont{part title}\insertpart\par
  \end{beamercolorbox}
}
\setbeamertemplate{section page}{
  \centering
  \begin{beamercolorbox}[sep=12pt,center]{part title}
    \usebeamerfont{section title}\insertsection\par
  \end{beamercolorbox}
}
\setbeamertemplate{subsection page}{
  \centering
  \begin{beamercolorbox}[sep=8pt,center]{part title}
    \usebeamerfont{subsection title}\insertsubsection\par
  \end{beamercolorbox}
}
\AtBeginPart{
  \frame{\partpage}
}
\AtBeginSection{
  \ifbibliography
  \else
    \frame{\sectionpage}
  \fi
}
\AtBeginSubsection{
  \frame{\subsectionpage}
}

\usepackage{amsmath,amssymb}
\usepackage{iftex}
\ifPDFTeX
  \usepackage[T1]{fontenc}
  \usepackage[utf8]{inputenc}
  \usepackage{textcomp} % provide euro and other symbols
\else % if luatex or xetex
  \usepackage{unicode-math}
  \defaultfontfeatures{Scale=MatchLowercase}
  \defaultfontfeatures[\rmfamily]{Ligatures=TeX,Scale=1}
\fi
\usepackage{lmodern}
\ifPDFTeX\else  
    % xetex/luatex font selection
\fi
% Use upquote if available, for straight quotes in verbatim environments
\IfFileExists{upquote.sty}{\usepackage{upquote}}{}
\IfFileExists{microtype.sty}{% use microtype if available
  \usepackage[]{microtype}
  \UseMicrotypeSet[protrusion]{basicmath} % disable protrusion for tt fonts
}{}
\usepackage{xcolor}
\newif\ifbibliography
\setlength{\emergencystretch}{3em} % prevent overfull lines
\setcounter{secnumdepth}{-\maxdimen} % remove section numbering


\providecommand{\tightlist}{%
  \setlength{\itemsep}{0pt}\setlength{\parskip}{0pt}}\usepackage{longtable,booktabs,array}
\usepackage{calc} % for calculating minipage widths
\usepackage{caption}
% Make caption package work with longtable
\makeatletter
\def\fnum@table{\tablename~\thetable}
\makeatother
\usepackage{graphicx}
\makeatletter
\def\maxwidth{\ifdim\Gin@nat@width>\linewidth\linewidth\else\Gin@nat@width\fi}
\def\maxheight{\ifdim\Gin@nat@height>\textheight\textheight\else\Gin@nat@height\fi}
\makeatother
% Scale images if necessary, so that they will not overflow the page
% margins by default, and it is still possible to overwrite the defaults
% using explicit options in \includegraphics[width, height, ...]{}
\setkeys{Gin}{width=\maxwidth,height=\maxheight,keepaspectratio}
% Set default figure placement to htbp
\makeatletter
\def\fps@figure{htbp}
\makeatother

\IfFileExists{plex-otf.sty}{
  %% Full TeXlive
  % \usepackage[%
  %   % math,
  %   RM={Scale=0.94},SS={Scale=0.94},SScon={Scale=0.94},TT={Scale=MatchLowercase,FakeStretch=0.9},DefaultFeatures={Ligatures=Common}
  % ]{plex-otf}
}{
  %% TinyTeX
  \usepackage{libertine}
}

%%% Load theme
% https://deic.uab.cat/~iblanes/beamer_gallery/
\IfFileExists{beamerthemegotham.sty}{
  %% Full TeXlive
  \usetheme{gotham}
  \gothamset{
    numbering=totalpagenumber,
    parttocframe default=off,
    sectiontocframe default=off,
    subsectiontocframe default=off,
  }
}{
  %% TinyTeX
  \usetheme{Madrid}
}
\makeatletter
\@ifpackageloaded{caption}{}{\usepackage{caption}}
\AtBeginDocument{%
\ifdefined\contentsname
  \renewcommand*\contentsname{Содержание}
\else
  \newcommand\contentsname{Содержание}
\fi
\ifdefined\listfigurename
  \renewcommand*\listfigurename{Список иллюстраций}
\else
  \newcommand\listfigurename{Список иллюстраций}
\fi
\ifdefined\listtablename
  \renewcommand*\listtablename{Список таблиц}
\else
  \newcommand\listtablename{Список таблиц}
\fi
\ifdefined\figurename
  \renewcommand*\figurename{Рисунок}
\else
  \newcommand\figurename{Рисунок}
\fi
\ifdefined\tablename
  \renewcommand*\tablename{Таблица}
\else
  \newcommand\tablename{Таблица}
\fi
}
\@ifpackageloaded{float}{}{\usepackage{float}}
\floatstyle{ruled}
\@ifundefined{c@chapter}{\newfloat{codelisting}{h}{lop}}{\newfloat{codelisting}{h}{lop}[chapter]}
\floatname{codelisting}{Список}
\newcommand*\listoflistings{\listof{codelisting}{Листинги}}
\makeatother
\makeatletter
\makeatother
\makeatletter
\@ifpackageloaded{caption}{}{\usepackage{caption}}
\@ifpackageloaded{subcaption}{}{\usepackage{subcaption}}
\makeatother

\ifLuaTeX
\usepackage[bidi=basic]{babel}
\else
\usepackage[bidi=default]{babel}
\fi
\babelprovide[main,import]{russian}
\babelprovide[import]{english}
% get rid of language-specific shorthands (see #6817):
\let\LanguageShortHands\languageshorthands
\def\languageshorthands#1{}
\ifLuaTeX
  \usepackage{selnolig}  % disable illegal ligatures
\fi
\usepackage{csquotes}
\usepackage{bookmark}

\IfFileExists{xurl.sty}{\usepackage{xurl}}{} % add URL line breaks if available
\urlstyle{same} % disable monospaced font for URLs
\hypersetup{
  pdftitle={Лабораторная работа 2},
  pdfauthor={Бондарь Татьяна Владимировна (я сменила свое имя)},
  pdflang={ru-RU},
  hidelinks,
  pdfcreator={LaTeX via pandoc}}


\title{Лабораторная работа 2}
\subtitle{Основы информационной безопасности}
\author{Бондарь Татьяна Владимировна (я сменила свое имя)}
\date{2026-03-05}

\begin{document}
\frame{\titlepage}

\renewcommand*\contentsname{Содержание}
\begin{frame}[allowframebreaks]
  \frametitle{Содержание}
  \tableofcontents[hideallsubsections]
\end{frame}

\section{Информация}\label{ux438ux43dux444ux43eux440ux43cux430ux446ux438ux44f}

\begin{frame}{Докладчик}
\phantomsection\label{ux434ux43eux43aux43bux430ux434ux447ux438ux43a}
\begin{columns}[c]
\begin{column}{0.7\textwidth}
Бондарь Татьяна Владимировна; студент группы НКАбд-01-24
\end{column}

\begin{column}{0.3\textwidth}
\end{column}
\end{columns}
\end{frame}

\section{Элементы
презентации}\label{ux44dux43bux435ux43cux435ux43dux442ux44b-ux43fux440ux435ux437ux435ux43dux442ux430ux446ux438ux438}

\begin{frame}{Цели и задачи}
\phantomsection\label{ux446ux435ux43bux438-ux438-ux437ux430ux434ux430ux447ux438}
Получение практических навыков работы в консоли с атрибутами файлов,
закрепление теоретических основ дискреционного разграничения доступа в
современных системах с открытым кодом на базе ОС Linux
\end{frame}

\section{Выполнение лабораторной
работы}\label{ux432ux44bux43fux43eux43bux43dux435ux43dux438ux435-ux43bux430ux431ux43eux440ux430ux442ux43eux440ux43dux43eux439-ux440ux430ux431ux43eux442ux44b}

\begin{frame}{1}
\phantomsection\label{section}
В операционной системе Rocky создаю нового пользователя guest через
учетную запись администратора (рис. 1).

\begin{figure}

{\centering \includegraphics[width=0.7\textwidth,height=\textheight]{image/1.PNG}

}

\caption{Добавление пользователя}

\end{figure}%
\end{frame}

\begin{frame}{2}
\phantomsection\label{section-1}
Сменяю пользователя в системе на только что созданного пользователя
guest (рис. 3).

\begin{figure}

{\centering \includegraphics[width=0.7\textwidth,height=\textheight]{image/3.PNG}

}

\caption{Вход через аккаунт пользователя}

\end{figure}%
\end{frame}

\begin{frame}{3}
\phantomsection\label{section-2}
В выводе команды groups информация только о названии группы, к которой
относится пользователь. В выводе команды id можно найти больше
информации: имя пользователя и имя группы, также коды имени пользователя
и группы (рис. 6)

\begin{figure}

{\centering \includegraphics[width=0.7\textwidth,height=\textheight]{image/6.PNG}

}

\caption{Информация о пользователе}

\end{figure}%
\end{frame}

\begin{frame}{4}
\phantomsection\label{section-3}
Имя пользователя в приглашении командной строкой совпадает с именем
пользователя, которое выводит команда whoami (рис. 7)

\begin{figure}

{\centering \includegraphics[width=0.7\textwidth,height=\textheight]{image/7.PNG}

}

\caption{Сравнение информации об имени пользователя}

\end{figure}%
\end{frame}

\begin{frame}{5}
\phantomsection\label{section-4}
Список поддиректорий директории home получилось получить с помощью
команды ls -l, если мы добавим опцию -a, то сможем увидеть еще и
директорию пользователя root. Права у директории:

root: drwxr-xr-x,

tvbondar и guest: drwx------ (рис. 9).
\end{frame}

\begin{frame}{6}
\phantomsection\label{section-5}
\begin{figure}

{\centering \includegraphics[width=0.7\textwidth,height=\textheight]{image/9.PNG}

}

\caption{Просмотр содержимого директории}

\end{figure}%
\end{frame}

\begin{frame}{7}
\phantomsection\label{section-6}
Создаю поддиректорию dir1 для домашней директории. Расширенные атрибуты
командой lsattr просмотреть у директории не удается, но атрибуты есть:
drwxr-xr-x, их удалось просмотреть с помощью команды ls -l (рис. 11).

\begin{figure}

{\centering \includegraphics[width=0.7\textwidth,height=\textheight]{image/11.PNG}

}

\caption{Создание поддиректории}

\end{figure}%
\end{frame}

\begin{frame}{8}
\phantomsection\label{section-7}
Снимаю атрибуты командой chmod 000 dir1, при проверке с помощью команды
ls -l видно, что теперь атрибуты действительно сняты (рис. 12).

\begin{figure}

{\centering \includegraphics[width=0.7\textwidth,height=\textheight]{image/12.PNG}

}

\caption{Снятие атрибутов с директории}

\end{figure}%
\end{frame}

\begin{frame}{9}
\phantomsection\label{section-8}
Попытка создать файл в директории dir1. Выдает ошибку: \enquote{Отказано
в доступе} (рис. 13).

\begin{figure}

{\centering \includegraphics[width=0.7\textwidth,height=\textheight]{image/13.PNG}

}

\caption{Попытка создания файла}

\end{figure}%
\end{frame}

\section{Выводы}\label{ux432ux44bux432ux43eux434ux44b}

\begin{frame}{Выводы}
Были получены практические навыки работы в консоли с атрибутами файлов,
закреплены теоретические основы дискреционного разграничения доступа в
современных системах с открытым кодом на базе ОС Linux.
\end{frame}




\end{document}
